\UTIResumenPagina
  {spanish}
  {Resumen Ejecutivo}
  {El presente trabajo de integración curricular aborda la actualización y modernización accesible del portal web oficial de la Facultad de Ingenierías de la Universidad Tecnológica Indoamérica (\url{https://facultad.indoamericatech.com/}), enfocado en garantizar el acceso autónomo a personas con discapacidad visual (ceguera total y baja visión). Partiendo de un diagnóstico técnico empírico realizado con Google Lighthouse v13.6.0 que evidenció deficiencias en contraste cromático, discontinuidades en la jerarquía de encabezados y discrepancias en nombres accesibles, se aplicó la metodología de diseño centrado en el humano (ISO 9241-210:2019) articulada con las directrices internacionales WCAG 2.2 Nivel AA. Se diseñaron e implementaron tres módulos prioritarios sobre una arquitectura moderna basada en Astro 7 y React 19: Oferta Académica y Mallas Curriculares, Trámites Estudiantiles y Preguntas Frecuentes, y un Asistente Conversacional Inteligente accesible con IA integrado mediante atajo de teclado (\texttt{Alt+A}) y anuncios dinámicos para el lector de pantalla NVDA mediante \texttt{aria-live="polite"}, complementado con una barra de herramientas de ajustes visuales. La validación empírica evidenció la resolución técnica de las barreras detectadas y una sustancial optimización en la experiencia de navegación autónoma, alcanzando una alta satisfacción de usuario y consolidando un entorno universitario plenamente inclusivo.}
  {Descriptores}
  {Accesibilidad web, Discapacidad visual, Diseño centrado en el usuario, Inteligencia artificial, WCAG 2.2.}
